\subsection{Models}
To analyze the relationship the yield curve with recessions, we employ a broad set of machine learning estimators. As a baseline estimator we use a simple Logistic regression.

\subsubsection{Regularization Approaches}

We employ regularization-driven estimation to identify the optimal linear predictors while mitigating the risk of overfitting through the application of $L_1$ and $L_2$ penalty terms, commonly referred to as selection and shrinkage, respectively. Specifically, we utilize a Ridge estimator, which relies on $L_2$ regularization to perform coefficient shrinkage, and an Elastic Net estimator, which incorporates a weighted combination of both $L_1$ and $L_2$ penalties. The Ridge regression estimator is defined by the following optimization objective, which introduces the $L_2$ penalty to the residual sum of squares:
\begin{equation}
\hat{\beta}_{Ridge} = \underset{\beta}{\text{argmin}} \left\{ \sum_{i=1}^{n} (y_i - x_i^T \beta)^2 + \lambda \sum_{j=1}^{p} \beta_j^2 \right\}
\end{equation}

The Elastic Net estimator bridges the gap between Ridge and Lasso by incorporating a mixing parameter $\alpha \in [0, 1]$, which balances the $L_1$ and $L_2$ penalties:
\begin{equation}
\hat{\beta}_{EN} = \underset{\beta}{\text{argmin}} \left\{ \frac{1}{2n} \|y - X\beta\|^2_2 + \lambda \left( \alpha \|\beta\|_1 + \frac{1-\alpha}{2} \|\beta\|^2_2 \right) \right\}
\end{equation}

To adapt these methods for classification, the optimization objective replaces the squared error used in regression with a log-likelihood loss function (log-loss) while maintaining the $L_1$ and $L_2$ penalty structures to constrain the model's coefficients.
\begin{equation}
\min_{w} \underbrace{\sum_{i=1}^{n} \log(1 + e^{-y_i x_i^T w})}_{\text{Classification Loss}} + \underbrace{\lambda \left[ \rho \|w\|_1 + \frac{1-\rho}{2} \|w\|_2^2 \right]}_{\text{Regularization Penalty}}
\end{equation}

Given that our feature set is constructed based on a priori economic rationale and supported by established academic literature, we hypothesize that automated feature selection may provide limited utility in this context. This assumption would be empirically supported if the Ridge estimator demonstrates predictive performance comparable or superior to that of the Elastic Net.

\subsubsection{Random Forests}
We include tree-based methods to strengthen our approach by relying on their non-parametric estimation properties. At their core, these methods estimate the conditional expectation $E[Y|X]$ by recursively partitioning the feature space into disjoint regions and assigning a constant value to each. For regression tasks, the model estimates the conditional average of the target variable within each leaf node, while in classification, the estimate represents the conditional probability of a class, which is effectively the conditional average of the class indicator variables. While a single tree approximates these averages through a simple local mean, ensemble architectures like Random Forest and Gradient Boosting enhance this estimation by either averaging across multiple independent partitions to reduce variance or sequentially refining the conditional mean to minimize a global loss function.

While both Random Forest and Boosting are ensemble methods based on decision trees, they differ fundamentally in how they combine learners. Random Forest uses bagging to build independent trees in parallel and averages their results to reduce variance, whereas Gradient Boosting and XGBoost employ boosting to build trees sequentially, with each new tree attempting to correct the errors of its predecessor. 

Specifically, XGBoost enhances the standard Gradient Boosting framework by incorporating $L_1$ and $L_2$ regularization directly into the objective function and utilizing a more efficient second-order Taylor expansion to optimize the loss. We anticipate that these boosting methods will capture complex interactions better than the Random Forest, provided the learning rate is sufficiently controlled to prevent overfitting.


Following ensemble logic and the Law of Large Numbers, whereby an aggregation of weakly correlated weak learners can yield a strong learner, the Random Forest prediction is defined simple average of $M$ independent trees:
\begin{equation}
\hat{y} = \frac{1}{M} \sum_{m=1}^{M} T_m(x)
\end{equation}


In contrast, Gradient Boosting produce an additive prediction where each step adds a small correction (weighted by the learning rate $\eta$):
\begin{equation}
\hat{y}_{new} = \hat{y}_{old} + \eta \cdot T_{next}(x)
\end{equation}


XGBoost further refines this by minimizing a regularized objective function that penalizes model complexity ($\Omega$):
\begin{equation}
\mathcal{L} = \sum \text{Loss}(y_i, \hat{y}_i) + \sum \Omega(T_m)
\end{equation}

\subsubsection{Metric Definitions}
\begin{itemize}
    \item \textbf{ROC\_AUC}: Measures the model's discriminative ability across all possible thresholds to separate signal from noise.
    \item \textbf{PR\_AUC}: Focuses on the precision of crisis predictions relative to recall to penalize false alarms.
    \item \textbf{Balanced Accuracy (BalAcc)}: The arithmetic mean of Recall and Specificity, ensuring the model is penalized for ignoring unlikely recession events.
    \item \textbf{F2 Score}: A weighted average of Precision and Recall that prioritizes capturing crises over perfect precision.
\end{itemize}

\subsection{Holdout Run}
\subsubsection{Short Term Holdout (Test Fraction 0.20)}
This period is likely dominated by post-GFC low-rate environments. The \textbf{Yield PCA (Logistic)} model achieved the highest performance (0.819 BalAcc, 0.213 F2) because PCA-cleared data allows simpler models to outperform complex ones like XGBoost. It successfully captured subtle curvature signals that were lost in the noisy \texttt{spread\_only} dataset. In contrast, the \textbf{Macro Only} model (0.530 BalAcc) indicates that real-sector data like unemployment was too slow for this short window.

\begin{table}[h!]
\centering
\caption{Holdout Sensitivity: Test Size 0.20}
\small
\begin{tabular}{@{}llllll@{}}
\toprule
\textbf{Feature Set} & \textbf{Model} & \textbf{ROC\_AUC} & \textbf{PR\_AUC} & \textbf{BalAcc} & \textbf{F2} \\ \midrule
Spreads Only & Logistic & 0.676 & 0.043 & 0.500 & 0.089 \\
& RandomForest & 0.613 & 0.036 & 0.353 & 0.000 \\
& GradBoost & 0.627 & 0.038 & 0.368 & 0.000 \\
& XGBoost & 0.544 & 0.030 & 0.730 & 0.154 \\ \midrule
Yield PCA & Logistic & 0.681 & 0.044 & 0.819 & 0.213 \\
& RandomForest & 0.667 & 0.042 & 0.500 & 0.089 \\
& GradBoost & 0.804 & 0.069 & 0.500 & 0.000 \\
& XGBoost & 0.627 & 0.038 & 0.417 & 0.000 \\ \midrule
Macro Only & Logistic & 0.505 & 0.029 & 0.530 & 0.097 \\
& RandomForest & 0.525 & 0.031 & 0.571 & 0.105 \\
& GradBoost & 0.389 & 0.024 & 0.480 & 0.000 \\
& XGBoost & 0.793 & 0.066 & 0.495 & 0.000 \\ \midrule
PCA + Macro & Logistic & 0.581 & 0.035 & 0.677 & 0.135 \\
& RandomForest & 0.677 & 0.044 & 0.621 & 0.118 \\
& GradBoost & 0.763 & 0.058 & 0.485 & 0.000 \\
& XGBoost & 0.783 & 0.063 & 0.629 & 0.152 \\ \midrule
PCA Macro Extras & Logistic & 0.512 & 0.036 & 0.585 & 0.128 \\
& RandomForest & 0.549 & 0.039 & 0.500 & 0.109 \\
& GradBoost & 0.360 & 0.036 & 0.457 & 0.000 \\
& XGBoost & 0.378 & 0.028 & 0.287 & 0.000 \\ \bottomrule
\end{tabular}
\end{table}


\subsubsection{Medium Term Holdout (Test Fraction 0.30)}
Covering the post-COVID recovery, this period saw noisy financial signals due to Quantitative Easing and Term Premium distortions. The \textbf{Yield PCA Gradient Boosting} model was the top performer (0.833 BalAcc, 0.928 ROC\_AUC) due to its ability to model complex nonlinear relations in noisy environments. However, the \textbf{PCA + Macro (Logistic)} model (0.787 BalAcc) is considered a safer, more stable choice for general usage as macro indicators rescued the model from failing financial signals.


\begin{table}[h!]
\centering
\caption{Holdout Sensitivity: Test Size 0.30}
\small
\begin{tabular}{@{}llllll@{}}
\toprule
\textbf{Feature Set} & \textbf{Model} & \textbf{ROC\_AUC} & \textbf{PR\_AUC} & \textbf{BalAcc} & \textbf{F2} \\ \midrule
Spreads Only & Logistic & 0.788 & 0.044 & 0.513 & 0.063 \\
& RandomForest & 0.631 & 0.025 & 0.611 & 0.078 \\
& GradBoost & 0.641 & 0.026 & 0.507 & 0.062 \\
& XGBoost & 0.559 & 0.021 & 0.489 & 0.056 \\ \midrule
Yield PCA & Logistic & 0.791 & 0.045 & 0.507 & 0.062 \\
& RandomForest & 0.637 & 0.026 & 0.549 & 0.068 \\
& GradBoost & 0.928 & 0.154 & 0.833 & 0.164 \\
& XGBoost & 0.765 & 0.040 & 0.735 & 0.110 \\ \midrule
Macro Only & Logistic & 0.683 & 0.030 & 0.730 & 0.110 \\
& RandomForest & 0.737 & 0.036 & 0.663 & 0.090 \\
& GradBoost & 0.567 & 0.022 & 0.487 & 0.000 \\
& XGBoost & 0.893 & 0.084 & 0.493 & 0.000 \\ \midrule
PCA + Macro & Logistic & 0.710 & 0.033 & 0.787 & 0.135 \\
& RandomForest & 0.800 & 0.047 & 0.657 & 0.088 \\
& GradBoost & 0.773 & 0.042 & 0.500 & 0.000 \\
& XGBoost & 0.840 & 0.057 & 0.660 & 0.139 \\ \midrule
PCA Macro Extras & Logistic & 0.698 & 0.038 & 0.758 & 0.143 \\
& RandomForest & 0.698 & 0.038 & 0.601 & 0.109 \\
& GradBoost & 0.560 & 0.027 & 0.355 & 0.000 \\
& XGBoost & 0.597 & 0.029 & 0.399 & 0.000 \\ \bottomrule
\end{tabular}
\end{table}

\subsubsection{Long Term Holdout (Test Fraction 0.40)}
This period includes major structural breaks like the 2008 GFC and COVID-19. The \textbf{Macro Only (Logistic)} model excelled (0.857 BalAcc) as instantly reacting unemployment and production data during shutdowns increased macro performance. While the Pure Yield PCA model failed due to the exogenous nature of the COVID shock, the \textbf{PCA + Macro (Logistic)} model (0.824 BalAcc) remains a robust candidate for prediction.

\begin{table}[h!]
\centering
\caption{Holdout Sensitivity: Test Size 0.40}
\small
\begin{tabular}{@{}llllll@{}}
\toprule
\textbf{Feature Set} & \textbf{Model} & \textbf{ROC\_AUC} & \textbf{PR\_AUC} & \textbf{BalAcc} & \textbf{F2} \\ \midrule
Spreads Only & Logistic & 0.552 & 0.062 & 0.494 & 0.146 \\
& RandomForest & 0.464 & 0.051 & 0.516 & 0.160 \\
& GradBoost & 0.492 & 0.052 & 0.426 & 0.000 \\
& XGBoost & 0.469 & 0.051 & 0.489 & 0.116 \\ \midrule
Yield PCA & Logistic & 0.562 & 0.063 & 0.527 & 0.167 \\
& RandomForest & 0.500 & 0.057 & 0.501 & 0.150 \\
& GradBoost & 0.482 & 0.062 & 0.560 & 0.195 \\
& XGBoost & 0.505 & 0.061 & 0.497 & 0.078 \\ \midrule
Macro Only & Logistic & 0.850 & 0.106 & 0.857 & 0.349 \\
& RandomForest & 0.701 & 0.055 & 0.503 & 0.133 \\
& GradBoost & 0.827 & 0.213 & 0.573 & 0.172 \\
& XGBoost & 0.642 & 0.057 & 0.500 & 0.000 \\ \midrule
PCA + Macro & Logistic & 0.762 & 0.066 & 0.824 & 0.303 \\
& RandomForest & 0.668 & 0.051 & 0.504 & 0.089 \\
& GradBoost & 0.915 & 0.377 & 0.731 & 0.294 \\
& XGBoost & 0.860 & 0.107 & 0.530 & 0.109 \\ \midrule
PCA Macro Extras & Logistic & 0.801 & 0.043 & 0.828 & 0.149 \\
& RandomForest & 0.768 & 0.037 & 0.807 & 0.135 \\
& GradBoost & 0.669 & 0.026 & 0.407 & 0.000 \\
& XGBoost & 0.681 & 0.027 & 0.370 & 0.000 \\ \bottomrule
\end{tabular}
\end{table}


\subsection{Scenario Analysis: COVID-19 and the Global Financial Crisis}
The analysis identifies a stark \textbf{Regime Dependence}:
\begin{itemize}
    \item \textbf{Financial Signal Victory}: During endogenous crises like 2008, Yield PCA factors successfully capture liquidity crunches with high precision.
    \item \textbf{Financial Signal Failure}: During exogenous shocks like COVID-19, the yield curve fails completely, and only real-sector macro data (Unemployment, Industrial Production) reacts fast enough.
\end{itemize}


\begin{itemize}
    \item \textbf{Regime-Switching Ensemble}: Implement a meta-model that weighs the Yield Curve model higher during monetary tightening and switches to Macro models during real-sector volatility.
    \item \textbf{High-Frequency Data}: Integrate weekly jobless claims or credit card spending to bridge the gap during "sudden stop" recessions where monthly data lags.
    \item \textbf{ACM Decomposition}: Isolate the "Risk-Neutral Slope" from the "Term Premium" to ensure robustness against QE distortions.
\end{itemize}

To assess regime dependence, models are evaluated separately on the COVID-19 recession (2019--2021) and the Global Financial Crisis (2007--2009). Figure \ref{fig:importance} report Variable Importance,  while Figures~\ref{fig:covid-pr} and~\ref{fig:gfc-pr} report precision--recall AUC.

\begin{figure}[H]
    \centering
    \includegraphics[width=0.85\textwidth]{figures/fig_importance_comparison_subplots.pdf}
    \caption{Top F2 scores during the COVID-19 recession.}
    \label{fig:importance}
\end{figure}

The feature importance analysis in demonstrates a clear divergence between the two economic regimes, validating the paper's core finding that recession signals are structurally dependent on the nature of the crisis. During the Global Financial Crisis (GFC), the slope factor (PC2) and the spread\_10y\_3m represent the most critical variables for the majority of models, particularly the Logistic and ElasticNet estimators. This dominance aligns with the characterization of the GFC as a financially driven recession where the erosion of the term premium and subsequent yield curve inversion served as primary barometers. In contrast, the COVID-19 scenario reveals a shift in importance toward PC1 (Level) and PC3 (Curvature), illustrating a structural break from traditional signals. The heightened relevance of the curvature factor (PC3) in the COVID-19 plot reflects the episodic spikes caused by intense policy intervention and stimulus programs that distorted the expectations-driven signals of the yield curve. These visual results provide empirical backing for the rolling PCA diagnostics, which show significant factor loading drift during these periods and underscore why models must be regime-aware to predict transitions effectively.

\begin{figure}[H]
    \centering
    \includegraphics[width=0.85\textwidth]{figures/fig_top_PR_AUC_COVID_test_2019_2021.pdf}
    \caption{Top PR-AUC scores during the COVID-19 recession.}
    \label{fig:covid-pr}
\end{figure}

\begin{figure}[H]
    \centering
    \includegraphics[width=0.85\textwidth]{figures/fig_top_PR_AUC_GFC_test_2007_2009.pdf}
    \caption{Top PR-AUC scores during the Global Financial Crisis.}
    \label{fig:gfc-pr}
\end{figure}

The scenario-based evaluation reveals significant regime dependence in the predictive power of yield curve factors. During the Global Financial Crisis (GFC) of 2007–2009, models leveraging principal component analysis (PCA) factors consistently outperformed those using raw maturity spreads, underscoring the informative value of slope and curvature for financially driven recessions. Specifically, the Ridge model with yield\_pca features achieved a Balanced Accuracy of 0.833 and a Recall of 0.722, significantly higher than the 0.694 Balanced Accuracy and 0.444 Recall obtained using raw spreads. This alignment with traditional financial theorywhere yield curve inversions effectively signal impending contractionsis further evidenced by the robust ROC\_AUC (0.938) and PR\_AUC (0.878) scores for the PCA-based Ridge model.

In contrast, the COVID-19 recession (2019–2021) represents a structural break where traditional yield signals were often lagging indicators of an abrupt, exogenous shock. Performance during this period was characterized by high Recall but lower Precision, reflecting the challenges of predicting a lockdown-induced downturn compared to an endogenous credit-cycle peak. While the ElasticNet and Logistic models with yield\_pca features both achieved a Recall of 1.000, their PR\_AUC values (0.174) were considerably lower than those seen during the GFC, with the exception of the Random Forest model which exhibited a perfect PR\_AUC of 1.000. This disparity reinforces the importance of regime-aware modeling, as unconventional monetary policies and term premium distortions in the post-2010 era have complicated the interpretation of latent yield factors.